% =====================================================================
% AI Academy -- National AI Center
% Software Engineering Final Project -- Team Report
% Topic 4 -- Async Research Assistant
% =====================================================================

\documentclass[11pt,a4paper]{article}

% ===== EARLY MACROS ==================================================
\usepackage[dvipsnames]{xcolor}
\definecolor{accent2}{HTML}{8B0000}
\definecolor{ph}{HTML}{6B7280}

% ===== CONFIGURATION =================================================
\newcommand{\teamname}{Allahverdi \& Isgandarli}
\newcommand{\topicchoice}{Topic 4 -- Async Research Assistant}
\newcommand{\member}[2]{\textbf{#1} (\texttt{#2})}
\newcommand{\reportdate}{18 September 2026}
\newcommand{\repolink}{\url{https://github.com/Yunis-Allahverdi/async-research-assistant}}

% ===== course meta ==================================================
\newcommand{\coursecode}{AI-ENG-110}
\newcommand{\coursename}{Software Engineering}
\newcommand{\institution}{AI Academy, National AI Center}
\newcommand{\semester}{Spring 2026}

% ===== PACKAGES =====================================================
\usepackage[margin=2.2cm,headheight=15pt]{geometry}
\usepackage[T1]{fontenc}
\usepackage[utf8]{inputenc}
\usepackage[final,protrusion=true,expansion=false]{microtype}
\usepackage{amsmath,amssymb}
\usepackage{graphicx}
\usepackage{booktabs}
\usepackage{array}
\usepackage{tabularx}
\usepackage{ragged2e}
\usepackage{enumitem}
\usepackage[parfill]{parskip}
\usepackage{listings}
\usepackage[most]{tcolorbox}
\usepackage{titlesec}
\usepackage{fancyhdr}
\usepackage{lastpage}
\usepackage[hidelinks]{hyperref}

\emergencystretch=3em
\hbadness=10000
\hfuzz=2pt

\hypersetup{
  colorlinks=true,
  linkcolor=NavyBlue,
  urlcolor=NavyBlue,
  citecolor=NavyBlue,
  pdftitle={\coursecode\ Final Project Report},
  pdfauthor={\teamname}
}

% ===== STYLING ======================================================
\definecolor{accent}{HTML}{0B3D91}
\definecolor{soft}{HTML}{F2F4F8}

\titleformat{\section}
  {\Large\bfseries\color{accent}}{\thesection.}{0.6em}{}
\titleformat{\subsection}
  {\large\bfseries\color{accent}}{\thesubsection}{0.5em}{}

\definecolor{codegreen}{rgb}{0,0.55,0}
\definecolor{codegray}{rgb}{0.4,0.4,0.4}
\definecolor{codepurple}{rgb}{0.55,0,0.78}
\definecolor{backcolour}{rgb}{0.97,0.97,0.97}

\lstdefinestyle{python}{
  language=Python,
  backgroundcolor=\color{backcolour},
  commentstyle=\color{codegreen},
  keywordstyle=\color{accent},
  numberstyle=\tiny\color{codegray},
  stringstyle=\color{codepurple},
  basicstyle=\ttfamily\footnotesize,
  breakatwhitespace=false,
  breaklines=true,
  keepspaces=true,
  numbers=left,
  numbersep=6pt,
  tabsize=2,
  frame=single,
  framerule=0.4pt,
  rulecolor=\color{black!30},
}
\lstdefinestyle{plain}{
  backgroundcolor=\color{backcolour},
  basicstyle=\ttfamily\footnotesize,
  breaklines=true,
  numbers=none,
  frame=single,
  framerule=0.4pt,
  rulecolor=\color{black!30},
}

\pagestyle{fancy}
\fancyhf{}
\fancyhead[L]{\small\textit{\coursecode\ -- Final Report -- \teamname}}
\fancyhead[R]{\small\textit{\semester}}
\fancyfoot[L]{\small\textit{\institution}}
\fancyfoot[R]{\small Page \thepage\ of \pageref{LastPage}}
\renewcommand{\headrulewidth}{0.4pt}
\renewcommand{\footrulewidth}{0.4pt}

% =====================================================================
\begin{document}

\begin{center}
{\small \textsc{\institution} \\[0.2em] \textsc{\coursecode\ -- \coursename\ -- \semester}}\\[1.2em]
{\Large\bfseries\color{accent} Final Project Report}\\[0.3em]
{\LARGE\bfseries \topicchoice}\\[0.6em]
\rule{0.85\textwidth}{0.6pt}
\end{center}

\vspace{0.6em}
\begin{tabularx}{\textwidth}{@{}lXlX@{}}
\textbf{Team:} & \teamname & \textbf{Date:} & \reportdate \\
\textbf{Repo:} & \repolink & \textbf{Tag:} & \texttt{v1.0-final} \\
\textbf{Members:} & \multicolumn{3}{X@{}}{%
  \member{Yunis Allahverdi}{@Yunis-Allahverdi},\quad
  \member{Fariz Isgandarli}{@farizisgenderli05-create}
}\\
\end{tabularx}
\vspace{0.4em}\hrule\vspace{0.6em}

% ===== EXECUTIVE SUMMARY =============================================
\section*{Executive Summary}
\addcontentsline{toc}{section}{Executive Summary}

We built an asynchronous research assistant: given a natural-language question, it queries
Wikipedia, arXiv, and a web-search provider concurrently, then uses a provider-agnostic LLM
to synthesize a single answer carrying inline \texttt{[N]} citations back to the retrieved
sources. The provided \texttt{ai/} package (fetchers plus citation synthesizer) was treated
as a frozen dependency; our work was the entire software-engineering layer around it ---
typed configuration, an \texttt{asyncio} orchestration layer, a filesystem cache, a retry /
timeout / logging service wrapper, input validation, a CLI, an offline test suite, and a
Docker image. The one number that proves the concurrency design works: on an identical
three-source workload our benchmark drops wall-clock time from \textbf{11.43\,s sequential to
3.64\,s concurrent (3.1$\times$)}. Our headline robustness story is that when \emph{all three}
live sources failed for reasons inside the frozen \texttt{ai/} layer, the system retried with
backoff, hit its per-source timeouts independently, and degraded to a clean ``no sources
found'' response rather than crashing. The most surprising lesson was that the hardest part
of an ``AI project'' was not the AI at all but everything that makes a thin API robust,
observable, and reproducible; if we started over we would push the retry policy to be
token-aware and add a real HTTP-header shim layer we are permitted to own.

% ===== TABLE OF CONTENTS ============================================
\tableofcontents
\vspace{0.6em}\hrule\vspace{0.4em}

% =====================================================================
\section{Project Overview}

\subsection{Problem framing \& scope}

The system answers a research question by gathering short, attributed excerpts from three
independent sources and letting a language model fuse them into one concise, cited answer.
The \emph{input} is a single question string; the \emph{outputs} are a synthesized answer
plus a numbered reference list, where every \texttt{[N]} marker in the prose maps to one
retrieved source. The only user-facing entry point is a command-line interface:
\texttt{python -m researcher ask "\dots"}, with \texttt{-{}-sources} to restrict the source
set and \texttt{-{}-no-cache} to bypass the cache. No HTTP server is required for this topic.

The provider stack is deliberately pluggable and selected entirely by environment variables:
the LLM is provider-agnostic (Anthropic / OpenAI / Gemini via \texttt{LLM\_PROVIDER}), and the
web-search backend is one of Tavily / Serper / DuckDuckGo via \texttt{WEB\_SEARCH\_PROVIDER}.
Wikipedia and arXiv need no key.

Relative to \texttt{TOPIC.md}, we implemented the full required SE layer and made two
deliberate scope calls. We \emph{tightened} the spec by making graceful degradation total:
not only does a single failing source not abort the call, but a call in which \emph{every}
source fails still returns a structured result instead of an error. We \emph{loosened} the
optional persistence choice: rather than PostgreSQL we used a filesystem-JSON cache, which is
persistent, inspectable, and adds no service to the container --- an intentional
simplicity-over-power trade-off for a pure key$\rightarrow$blob workload.

\subsection{Provider choice and the AI module contract}

The LLM is selected at runtime through the provided factory; for development the default is
\texttt{LLM\_PROVIDER=anthropic} with model \texttt{claude-sonnet-4-6}, and the web backend
default is \texttt{WEB\_SEARCH\_PROVIDER=duckduckgo} (no key, convenient for grading). Our
code calls exactly four functions of the frozen package --- \texttt{ai.fetch\_wikipedia},
\texttt{ai.fetch\_arxiv}, \texttt{ai.fetch\_web}, and \texttt{ai.synthesize} --- and the
schemas \texttt{Source}, \texttt{Citation}, \texttt{AnswerWithCitations}.
\textbf{We did not modify the public interface of the provided \texttt{ai/} package in any
way.} Against malformed model output we rely on two guards: the provided synthesizer already
drops out-of-range / hallucinated citation indices, and our service layer validates that a
call returned a usable, non-empty result before caching or returning it, treating provider
faults as the transient \texttt{ProviderError} class described in \S\ref{sec:robustness}.

% =====================================================================
\section{Architecture}\label{sec:arch}

\subsection{Module map}

The system is a strict downward-dependency pipeline. The CLI drives the core researcher, which
drives the concurrency orchestrator, which fans out to the service wrapper, which is the
\emph{only} module that touches the frozen \texttt{ai/} package, the cache, and the config.

\begin{lstlisting}[style=plain]
        CLI  (src/cli.py, researcher/__main__.py)
         |  ask "question" [--sources ...] [--no-cache]
         v
   core/researcher.py        validate -> orchestrate -> synthesize -> ResearchResult
         |
         v
 concurrency/orchestrator.py  asyncio.gather(return_exceptions=True)
         |                    + Semaphore(max_concurrency) + per-source asyncio.timeout
         |  fetch_source(origin, query)  x3, one shared httpx.AsyncClient
         v
 services/ai_service.py       retry+backoff | timeout | logging | cache lookup/store
     |        |         |
     v        v         v
 storage/  ===AI BOUNDARY===   config.py
 cache_store   ai/ (FROZEN)    Settings (env-driven)
 (JSON+TTL)    fetch_* / synthesize
\end{lstlisting}

\textbf{Module-by-module.}
\texttt{config.py} reads \texttt{.env} into a typed \texttt{Settings}.
\texttt{models.py} holds our SE-layer schemas.
\texttt{storage/cache\_store.py} is a TTL-aware filesystem JSON cache.
\texttt{services/ai\_service.py} wraps every \texttt{ai.*} call with retry, timeout, logging,
and caching --- the single crossing of the AI boundary.
\texttt{concurrency/orchestrator.py} runs the three fetches concurrently with bounded
parallelism and graceful degradation.
\texttt{core/researcher.py} validates input and assembles the final \texttt{ResearchResult}.
\texttt{cli.py} parses arguments and renders the answer.
Only the arrows into \texttt{ai\_service} cross the AI boundary; swapping the LLM provider
(Anthropic~$\rightarrow$~OpenAI) changes \emph{zero} of our files --- it is one environment
variable, because the provided factory hides the concrete provider behind an ABC.

\subsection{OOP design \& key abstractions}

Our own defensible abstraction is \emph{composition via a dispatch registry} in the service
layer, rather than inheritance. The three source fetchers share one call shape, so instead of
an \texttt{if/elif} ladder or a subclass per source we map origin~$\rightarrow$~coroutine and
inject the shared client and settings:

\begin{lstlisting}[style=python]
_FETCHERS = {
    "wikipedia": fetch_wikipedia,
    "arxiv": fetch_arxiv,
    "web": fetch_web,
}

async def fetch_source(origin, query, *, client=None, settings, use_cache=True):
    fetcher = _FETCHERS.get(origin)
    if fetcher is None:
        raise ValueError(f"unknown origin {origin!r}")
    ...
    return await _fetch_with_retry(fetcher, query, client=client, settings=settings)
\end{lstlisting}

\textbf{Justification.} We chose composition because the fetchers are stateless functions
with an identical signature --- there is no shared state or behaviour to inherit, so a class
hierarchy would add ceremony without value. A registry keeps ``add a source'' to one line and
keeps retry/timeout/caching in exactly one place. Where the problem genuinely called for an
open provider set with per-provider state (API keys, endpoints), the \emph{provided} code
correctly used an ABC (\texttt{WebSearchProvider} with Tavily/Serper/DuckDuckGo subclasses);
we deliberately did not duplicate that pattern where plain functions suffice.

\subsection{Data flow across module boundaries}

We defined three SE-layer models in \texttt{src/models.py}: \texttt{SourceResult} (per-source
outcome: origin, sources, \texttt{ok}, error, \texttt{elapsed}), \texttt{ResearchResult}
(question, answer, per-source results, \texttt{degraded} flag, total elapsed, \texttt{from\_cache}),
and \texttt{CacheEntry} (key, sources, \texttt{created\_at}). We reused the provided
\texttt{Source} / \texttt{AnswerWithCitations} verbatim wherever the data was exactly what
\texttt{ai/} produced, and wrapped them only when we needed to attach SE concerns the provided
schema does not carry --- timing and success/failure for the orchestrator, cache metadata for
the store. \textbf{No naked dictionaries cross module boundaries; every inter-module value is
a typed pydantic model.}

% =====================================================================
\section{Concurrency \& Performance}\label{sec:concurrency}

\subsection{Why this concurrency model}

The workload is I/O-bound --- three independent network round-trips that spend nearly all
their wall-clock time waiting on remote servers, not on CPU. \texttt{asyncio} is therefore the
right primitive: a single thread can keep all three requests in flight, and there is no CPU
work to parallelize that would justify \texttt{ProcessPoolExecutor}, nor any blocking SDK that
would force a \texttt{ThreadPoolExecutor}. Threads would also work for three requests but cost
one OS thread per source and a shared connection pool is cleaner under one event loop. The
bottleneck the design targets is \emph{total network latency}: sequentially it is the sum of
the three sources; concurrently it collapses to the slowest single source. We bound
parallelism with \texttt{asyncio.Semaphore(max\_concurrency)} (default 5, configurable via
\texttt{MAX\_CONCURRENCY}); at 1 the design degenerates to sequential, and at a very high value
we would risk a provider rate-limit (HTTP 429), so a small bound respects politeness while
still overlapping our three-way fan-out. The concurrent design starts winning at $N\ge2$
sources.

\subsection{Sequential-vs-concurrent benchmark}

\texttt{scripts/benchmark.py} runs the same query twice --- once awaiting each source in turn,
once through the orchestrator --- always bypassing the cache so both do identical work.
Measured on a Windows 11 laptop, Python 3.12.4, caches cleared:

\begin{center}\small
\begin{tabular}{lcccc}
\toprule
\textbf{Workload} & $N$ & \textbf{Sequential (s)} & \textbf{Concurrent (s)} & \textbf{Speedup} \\
\midrule
1 query, 3 sources (wiki, arxiv, web) & 3 & 11.43 & 3.64 & 3.1$\times$ \\
\bottomrule
\end{tabular}
\end{center}

Command:
\begin{lstlisting}[style=plain]
python -m scripts.benchmark "What is photosynthesis?" --sources wiki,arxiv,web
\end{lstlisting}

\textbf{Honest interpretation.} We reached almost exactly the theoretical $3\times$ for a
three-source fan-out, which is expected because the three latencies were comparable and the
work is pure waiting. An important caveat: during measurement all three live endpoints were
failing (see \S\ref{sec:robustness}), so the numbers time the concurrent-vs-sequential
handling of the retry-and-timeout path rather than of successful downloads --- but the
comparison is still valid, because sequential pays each source's retry/timeout budget in turn
(sum) while concurrent pays them in parallel (max). The new bottleneck after parallelizing the
fetch is the \emph{LLM synthesis step}, which is a single call that runs after all fetches
complete and is therefore serial tail latency; pushing the curve further would mean streaming
the synthesis or overlapping it with late-arriving sources.

% =====================================================================
\section{Robustness \& Error Handling}\label{sec:robustness}

\subsection{Wrapping the AI module}

Every \texttt{ai.*} call goes through \texttt{ai\_service}, which adds four things the thin
provided layer lacks. \textbf{Retries:} a \texttt{tenacity} policy retries \emph{only}
\texttt{ProviderError} --- the transient class (network errors, 5xx, 429) --- for up to
\texttt{MAX\_RETRIES} (default 3) attempts with \texttt{wait\_exponential} backoff.
Non-transient failures such as a \texttt{ValueError} from empty input are \emph{not} retried,
because they cannot succeed on repeat and retrying them only wastes the budget.
\textbf{Timeouts:} each attempt is wrapped in \texttt{asyncio.timeout(source\_timeout)}
(default 10\,s) as a hard per-call ceiling. \textbf{Structured logging:} the standard
\texttt{logging} module (never \texttt{print}) records origin, truncated query, result count,
and elapsed time at INFO, full payloads at DEBUG, and each retry at WARNING; the level is set
once from \texttt{LOG\_LEVEL}. \textbf{Validation:} input is checked before the call (see
\S\ref{sec:validation}) and the result is checked for usability before it is cached or
returned. Walking one case end-to-end: a provider 429 is raised by \texttt{ai/} as a
\texttt{ProviderError}, caught by the retry predicate, logged at WARNING, waited-out with
exponential backoff, and retried; after the attempt budget is exhausted the original error is
re-raised (\texttt{reraise=True}) and surfaces to the orchestrator, which records it as a
failed \texttt{SourceResult} without cancelling the other sources.

\subsection{Failure-mode analysis (one concrete incident)}

\textbf{Incident: all three live sources unreachable.} We did not have to inject this failure
--- it occurred naturally, which made it a stronger test. Running
\texttt{python -m researcher ask "How does CRISPR-Cas9 work?"} against the live network, every
source failed for a reason inside the frozen \texttt{ai/} layer: arXiv returned HTTP 301 (the
fetcher issues an \texttt{http://} request and does not follow the redirect to
\texttt{https://}), Wikipedia returned HTTP 403 (the request carries no \texttt{User-Agent}
header), and the pinned \texttt{duckduckgo-search} dependency had been renamed to \texttt{ddgs}
upstream and now returns an empty list.

\emph{(i) How it surfaced:} a normal live query. \emph{(ii) What the system did:} each source
was retried three times with backoff, each hit its timeout independently, and the orchestrator
collected three failed \texttt{SourceResult}s via \texttt{return\_exceptions=True}.
\emph{(iii) What the user saw:} a clean message --- ``No sources could be found for this
question\dots'' followed by ``\emph{Note: wikipedia, arxiv, web could not be reached, so this
answer may be incomplete}'' --- and no stack trace. \emph{(iv) What the logs showed:} an
ordered WARNING trail of \texttt{retrying (1/3)}, \texttt{(2/3)}, then per-source
\texttt{FAILED} lines naming the exact HTTP status, which was enough to diagnose all three
root causes without re-running the code. We originally expected at least DuckDuckGo (keyless)
to succeed; the log is what revealed the upstream package rename. Because \texttt{ai/} is
contractually frozen we could not patch the redirect/header/rename at the source, so the
correct engineering response was exactly what the system already did --- degrade cleanly and
report --- which validated the robustness path against real endpoints.

\subsection{Input validation}\label{sec:validation}

The single entry point is the CLI. \texttt{validate\_question} strips the question and rejects
it if it is shorter than \texttt{MIN\_QUESTION\_LEN} (default 3) or longer than
\texttt{MAX\_QUESTION\_LEN} (default 500), raising a \texttt{ValueError} that the CLI surfaces
as a one-line message rather than a traceback. The \texttt{-{}-sources} flag is validated
against the allowed set \texttt{\{wiki, arxiv, web\}}; an unknown origin raises a clear error.
Environment values are validated by \texttt{pydantic-settings} at load time (wrong types fail
fast). A 50{,}000-character query is rejected by the length bound before any network call, so
an adversarial oversized input cannot reach a provider; cache keys are hashed, so a query can
never escape the storage directory via crafted path characters.

\subsection{Rate limit \& politeness}

Politeness has two layers: the \texttt{Semaphore} bounds concurrent in-flight requests so we
never burst more than \texttt{MAX\_CONCURRENCY} (default 5) at once, and the exponential-backoff
retry acts as a sleep-on-429 strategy --- a provider 429 is a \texttt{ProviderError}, so the
backoff naturally spaces out subsequent attempts. The highest burst our test suite produces is
three simultaneous fetches (one per source), well under the bound. We did not observe a real
429 in development because the live endpoints failed earlier in the request path; the strategy
is therefore exercised by unit tests that inject \texttt{ProviderError} rather than by a
production 429.

% =====================================================================
\section{Storage \& Caching}\label{sec:storage}

\subsection{Design}

We did not use SQLite/PostgreSQL. The workload is a pure key$\rightarrow$value blob store
(\texttt{(origin, query)} $\rightarrow$ list of \texttt{Source}s), for which a relational
schema would be over-engineering and would add a service to the container. Instead each entry
is one JSON file under \texttt{CACHE\_DIR} (default \texttt{.cache/}), named by a stable key.
The key is the source-of-truth of identity; the stored \texttt{Source} list and
\texttt{created\_at} timestamp are the cached payload.

\begin{lstlisting}[style=python]
def make_key(origin, query):
    canonical = f"{origin.lower().strip()}|{query.lower().strip()}"
    return f"{origin}_{hashlib.sha256(canonical.encode()).hexdigest()[:16]}"
\end{lstlisting}

The key is \emph{canonicalised} (lower-cased, stripped) so ``\texttt{WHAT IS X?}'' and
``\texttt{what is x}'' collide onto one entry, then hashed so arbitrary query characters
cannot break the filename or escape the directory.

\subsection{Caching policy}

We cache each \texttt{(origin, query)} result for \texttt{CACHE\_TTL} seconds (default 86{,}400
= 24\,h). On read, the entry's age is compared to the TTL and anything older is treated as a
miss and refetched --- this is time-based invalidation. \texttt{ai\_service} checks the cache
before every fetch and writes non-empty results after; \texttt{-{}-no-cache} bypasses both ends.
In a demo run that issues the same query twice, the first call is a miss and the second is a
\textbf{100\% cache hit} (the underlying fetcher is invoked exactly once, verified by a unit
test that counts fetcher calls). The value most likely to go stale is an LLM-shaped result if
the prompt template changes; in production we would add a prompt-version component to the key
so a changed prompt misses the old entries automatically.

% =====================================================================
\section{Testing}\label{sec:testing}

\subsection{Strategy \& coverage}

Every test runs \textbf{offline}: each \texttt{ai.*} call is replaced with a fake via
\texttt{monkeypatch}, so the suite needs no network and no API keys and is safe in CI. Async
code is driven by \texttt{pytest-asyncio} (\texttt{asyncio\_mode = strict}). We chose not to
chase the last few lines --- three uncovered lines in \texttt{config.py} are the
\texttt{\_\_main\_\_} print block, and the two in the orchestrator are a defensive branch that
only fires on a malformed internal result --- because covering them would test framework
plumbing, not our logic.

\begin{center}\small
\begin{tabular}{lc}
\toprule
\textbf{Suite} & \textbf{Coverage} \\
\midrule
\texttt{src/} (our SE layer) & 98\,\% (249 stmts, 6 missed) \\
Provided \texttt{ai/} smoke tests & passing \\
\bottomrule
\end{tabular}
\end{center}

The full suite is \textbf{51 tests passing} in $\sim$2\,s, and the provided
\texttt{tests/test\_ai\_smoke.py} is unmodified and passes.
A \texttt{mypy src} run reports 6 warnings: one is in the frozen \texttt{ai/} package (which we
may not edit), and five are benign false positives in \texttt{ai\_service.py} --- mypy cannot
prove the three registered fetchers share a signature, nor that the retry loop
(\texttt{reraise=True}) always assigns before use. All are documented; none affects runtime,
as the 51 passing tests confirm.

\subsection{Notable tests}

\textbf{(i) Happy path (end-to-end).} \texttt{test\_fetch\_source\_happy\_path} mocks a fetcher
and asserts a well-formed \texttt{Source} flows back through the cache-write path.
\textbf{(ii) Concurrency.} \texttt{test\_orchestrator} verifies that when one of three
gathered fetches raises, \texttt{return\_exceptions=True} still yields three
\texttt{SourceResult}s --- two \texttt{ok}, one failed --- and no exception propagates.
\textbf{(iii) Error path.} \texttt{test\_fetch\_source\_retries\_then\_succeeds} and
\texttt{\dots\_reraises\_after\_exhaustion} pin the retry contract: a \texttt{ProviderError}
that clears on the third attempt succeeds, while a permanent one is re-raised after exactly
\texttt{MAX\_RETRIES} attempts. \texttt{\dots\_times\_out} confirms a slow fetcher raises
\texttt{TimeoutError} and is \emph{not} retried. The test we wish we had time for is a
load-style test asserting the semaphore never lets more than \texttt{MAX\_CONCURRENCY}
coroutines run at once under a large fan-out.

% =====================================================================
\section{Deployment \& Reproducibility}\label{sec:deploy}

\subsection{Docker image}

\begin{lstlisting}[style=plain]
docker build -t research-assistant .
docker run --rm research-assistant        # runs: python demo_ai.py --offline
\end{lstlisting}

It is a \textbf{single-stage} build on \texttt{python:3.12-slim}. The image is \textbf{76.5\,MB
content} (324\,MB on-disk including the shared base layers), comparable to a fresh
\texttt{python:3.12-slim} ($\sim$120\,MB) because our only added weight is the pinned Python
dependencies and the small source tree; a \texttt{.dockerignore} keeps the virtualenv, caches,
git history, and \texttt{.env} out of the image. Dependencies are installed before the source
is copied so the dependency layer is cached across code changes. The default command runs the
offline demo end-to-end with no keys or network; the live CLI can be run by passing env vars
with \texttt{-e} / \texttt{-{}-env-file}.

\subsection{Configuration}

Every variable the app reads, and the effect if it is missing:
\texttt{LLM\_PROVIDER} (default \texttt{anthropic}); the provider's API key
(\texttt{ANTHROPIC\_API\_KEY} etc.) --- if missing, live synthesis raises a clear
\texttt{ProviderError}, but the offline demo and tests are unaffected;
\texttt{WEB\_SEARCH\_PROVIDER} (default \texttt{duckduckgo}); \texttt{TAVILY\_API\_KEY} /
\texttt{SERPER\_API\_KEY} (needed only for those backends); \texttt{CACHE\_DIR} (default
\texttt{.cache}); \texttt{CACHE\_TTL} (default 86400); \texttt{SOURCE\_TIMEOUT} (default 10);
\texttt{MAX\_CONCURRENCY} (default 5); \texttt{MAX\_RETRIES} / \texttt{RETRY\_BACKOFF}
(defaults 3 / 1.0); \texttt{MIN\_QUESTION\_LEN} / \texttt{MAX\_QUESTION\_LEN} (defaults 3 /
500); \texttt{LOG\_LEVEL} (default \texttt{INFO}). Every non-key variable has a safe default,
so the app runs with an empty \texttt{.env}.

% =====================================================================
\section{Limitations \& Future Work}\label{sec:limits}

\begin{itemize}[leftmargin=1.4em,itemsep=3pt]
  \item \textbf{Live retrieval depends on frozen code we cannot patch.} arXiv (301 redirect),
  Wikipedia (403 / missing User-Agent), and the renamed \texttt{duckduckgo-search} package all
  fail inside \texttt{ai/}; we validate the pipeline via the offline demo and mocked tests. A
  permitted fix would be a thin request-shim layer of our own that adds headers and follows
  redirects \emph{before} handing off, without editing \texttt{ai/}.
  \item \textbf{The filesystem cache has no concurrent-writer locking} and grows one file per
  unique query; a production version would add locking or move to SQLite/Redis and an eviction
  policy.
  \item \textbf{Synthesis is serial tail latency} --- the single LLM call runs after all
  fetches, so it dominates once retrieval is parallel; streaming would help.
  \item \textbf{Rate limiting is request-count only}, not token-aware; a busy day with a real
  provider could still approach a tokens-per-minute limit the semaphore does not see.
  \item \textbf{Two-person scope} meant we prioritized the required layer over optional bonuses
  (multi-provider failover, cost telemetry, CI); these are the natural next week of work.
\end{itemize}

% =====================================================================
\section{Tools \& Acknowledgements}\label{sec:tools}

We used AI coding assistants as collaborators, and disclose this specifically rather than
generically. \textbf{Claude} was used to scaffold the SE layer module by module and to explain
each concept (\texttt{asyncio}, \texttt{tenacity} retry semantics, \texttt{pytest} mocking,
Docker) as we built; every module was then reviewed, run, and unit-tested by its owner, and we
can each defend our own files. Concrete example: the \texttt{ai\_service} retry policy was
drafted with Claude, but we tuned it after reading the live logs --- narrowing the retry
predicate to \texttt{ProviderError} only, so that non-transient \texttt{ValueError}s are not
uselessly retried --- and we added the per-attempt \texttt{asyncio.timeout} after observing
that a hung source would otherwise block within the retry budget. The provided \texttt{ai/}
package, its schemas, and \texttt{demo\_ai.py} are course-supplied and were left unmodified.

% =====================================================================
\section*{References}
\addcontentsline{toc}{section}{References}

\begin{enumerate}[label={[\arabic*]},leftmargin=2.4em,itemsep=2pt]
  \item Anthropic API documentation. \url{https://docs.anthropic.com/}
  \item \texttt{tenacity} --- retrying library. \url{https://tenacity.readthedocs.io/}
  \item \texttt{pydantic-settings} documentation. \url{https://docs.pydantic.dev/latest/concepts/pydantic_settings/}
  \item \texttt{httpx} --- async HTTP client. \url{https://www.python-httpx.org/}
  \item \texttt{pytest-asyncio} documentation. \url{https://pytest-asyncio.readthedocs.io/}
  \item Python \texttt{asyncio} --- \texttt{gather}, \texttt{timeout}, \texttt{Semaphore}. \url{https://docs.python.org/3/library/asyncio.html}
\end{enumerate}

\appendix
% =====================================================================
\section{Appendix --- Reproducing the benchmark}
Anyone with the repo and a Python 3.11+ environment can reproduce the \S\ref{sec:concurrency}
table:

\begin{lstlisting}[style=plain]
git clone https://github.com/Yunis-Allahverdi/async-research-assistant
cd async-research-assistant
python -m venv venv && source venv/bin/activate   # Windows: venv\Scripts\Activate.ps1
python -m pip install -r requirements-ai.txt -r requirements.txt
python -m scripts.benchmark "What is photosynthesis?" --sources wiki,arxiv,web
\end{lstlisting}

\vspace{1em}\hrule\vspace{0.6em}
\noindent\small\textit{End of report --- thank you for reading.}

\end{document}% =====================================================================
% AI Academy -- National AI Center
% Software Engineering Final Project -- Team Report
% Topic 4 -- Async Research Assistant
% =====================================================================

\documentclass[11pt,a4paper]{article}

% ===== EARLY MACROS ==================================================
\usepackage[dvipsnames]{xcolor}
\definecolor{accent2}{HTML}{8B0000}
\definecolor{ph}{HTML}{6B7280}

% ===== CONFIGURATION =================================================
\newcommand{\teamname}{Allahverdi \& Isgandarli}
\newcommand{\topicchoice}{Topic 4 -- Async Research Assistant}
\newcommand{\member}[2]{\textbf{#1} (\texttt{#2})}
\newcommand{\reportdate}{18 September 2026}
\newcommand{\repolink}{\url{https://github.com/Yunis-Allahverdi/async-research-assistant}}

% ===== course meta ==================================================
\newcommand{\coursecode}{AI-ENG-110}
\newcommand{\coursename}{Software Engineering}
\newcommand{\institution}{AI Academy, National AI Center}
\newcommand{\semester}{Spring 2026}

% ===== PACKAGES =====================================================
\usepackage[margin=2.2cm,headheight=15pt]{geometry}
\usepackage[T1]{fontenc}
\usepackage[utf8]{inputenc}
\usepackage[final,protrusion=true,expansion=false]{microtype}
\usepackage{amsmath,amssymb}
\usepackage{graphicx}
\usepackage{booktabs}
\usepackage{array}
\usepackage{tabularx}
\usepackage{ragged2e}
\usepackage{enumitem}
\usepackage[parfill]{parskip}
\usepackage{listings}
\usepackage[most]{tcolorbox}
\usepackage{titlesec}
\usepackage{fancyhdr}
\usepackage{lastpage}
\usepackage[hidelinks]{hyperref}

\emergencystretch=3em
\hbadness=10000
\hfuzz=2pt

\hypersetup{
  colorlinks=true,
  linkcolor=NavyBlue,
  urlcolor=NavyBlue,
  citecolor=NavyBlue,
  pdftitle={\coursecode\ Final Project Report},
  pdfauthor={\teamname}
}

% ===== STYLING ======================================================
\definecolor{accent}{HTML}{0B3D91}
\definecolor{soft}{HTML}{F2F4F8}

\titleformat{\section}
  {\Large\bfseries\color{accent}}{\thesection.}{0.6em}{}
\titleformat{\subsection}
  {\large\bfseries\color{accent}}{\thesubsection}{0.5em}{}

\definecolor{codegreen}{rgb}{0,0.55,0}
\definecolor{codegray}{rgb}{0.4,0.4,0.4}
\definecolor{codepurple}{rgb}{0.55,0,0.78}
\definecolor{backcolour}{rgb}{0.97,0.97,0.97}

\lstdefinestyle{python}{
  language=Python,
  backgroundcolor=\color{backcolour},
  commentstyle=\color{codegreen},
  keywordstyle=\color{accent},
  numberstyle=\tiny\color{codegray},
  stringstyle=\color{codepurple},
  basicstyle=\ttfamily\footnotesize,
  breakatwhitespace=false,
  breaklines=true,
  keepspaces=true,
  numbers=left,
  numbersep=6pt,
  tabsize=2,
  frame=single,
  framerule=0.4pt,
  rulecolor=\color{black!30},
}
\lstdefinestyle{plain}{
  backgroundcolor=\color{backcolour},
  basicstyle=\ttfamily\footnotesize,
  breaklines=true,
  numbers=none,
  frame=single,
  framerule=0.4pt,
  rulecolor=\color{black!30},
}

\pagestyle{fancy}
\fancyhf{}
\fancyhead[L]{\small\textit{\coursecode\ -- Final Report -- \teamname}}
\fancyhead[R]{\small\textit{\semester}}
\fancyfoot[L]{\small\textit{\institution}}
\fancyfoot[R]{\small Page \thepage\ of \pageref{LastPage}}
\renewcommand{\headrulewidth}{0.4pt}
\renewcommand{\footrulewidth}{0.4pt}

% =====================================================================
\begin{document}

\begin{center}
{\small \textsc{\institution} \\[0.2em] \textsc{\coursecode\ -- \coursename\ -- \semester}}\\[1.2em]
{\Large\bfseries\color{accent} Final Project Report}\\[0.3em]
{\LARGE\bfseries \topicchoice}\\[0.6em]
\rule{0.85\textwidth}{0.6pt}
\end{center}

\vspace{0.6em}
\begin{tabularx}{\textwidth}{@{}lXlX@{}}
\textbf{Team:} & \teamname & \textbf{Date:} & \reportdate \\
\textbf{Repo:} & \repolink & \textbf{Tag:} & \texttt{v1.0-final} \\
\textbf{Members:} & \multicolumn{3}{X@{}}{%
  \member{Yunis Allahverdi}{@Yunis-Allahverdi},\quad
  \member{Fariz Isgandarli}{@farizisgenderli05-create}
}\\
\end{tabularx}
\vspace{0.4em}\hrule\vspace{0.6em}

% ===== EXECUTIVE SUMMARY =============================================
\section*{Executive Summary}
\addcontentsline{toc}{section}{Executive Summary}

We built an asynchronous research assistant: given a natural-language question, it queries
Wikipedia, arXiv, and a web-search provider concurrently, then uses a provider-agnostic LLM
to synthesize a single answer carrying inline \texttt{[N]} citations back to the retrieved
sources. The provided \texttt{ai/} package (fetchers plus citation synthesizer) was treated
as a frozen dependency; our work was the entire software-engineering layer around it ---
typed configuration, an \texttt{asyncio} orchestration layer, a filesystem cache, a retry /
timeout / logging service wrapper, input validation, a CLI, an offline test suite, and a
Docker image. The one number that proves the concurrency design works: on an identical
three-source workload our benchmark drops wall-clock time from \textbf{11.43\,s sequential to
3.64\,s concurrent (3.1$\times$)}. Our headline robustness story is that when \emph{all three}
live sources failed for reasons inside the frozen \texttt{ai/} layer, the system retried with
backoff, hit its per-source timeouts independently, and degraded to a clean ``no sources
found'' response rather than crashing. The most surprising lesson was that the hardest part
of an ``AI project'' was not the AI at all but everything that makes a thin API robust,
observable, and reproducible; if we started over we would push the retry policy to be
token-aware and add a real HTTP-header shim layer we are permitted to own.

% ===== TABLE OF CONTENTS ============================================
\tableofcontents
\vspace{0.6em}\hrule\vspace{0.4em}

% =====================================================================
\section{Project Overview}

\subsection{Problem framing \& scope}

The system answers a research question by gathering short, attributed excerpts from three
independent sources and letting a language model fuse them into one concise, cited answer.
The \emph{input} is a single question string; the \emph{outputs} are a synthesized answer
plus a numbered reference list, where every \texttt{[N]} marker in the prose maps to one
retrieved source. The only user-facing entry point is a command-line interface:
\texttt{python -m researcher ask "\dots"}, with \texttt{-{}-sources} to restrict the source
set and \texttt{-{}-no-cache} to bypass the cache. No HTTP server is required for this topic.

The provider stack is deliberately pluggable and selected entirely by environment variables:
the LLM is provider-agnostic (Anthropic / OpenAI / Gemini via \texttt{LLM\_PROVIDER}), and the
web-search backend is one of Tavily / Serper / DuckDuckGo via \texttt{WEB\_SEARCH\_PROVIDER}.
Wikipedia and arXiv need no key.

Relative to \texttt{TOPIC.md}, we implemented the full required SE layer and made two
deliberate scope calls. We \emph{tightened} the spec by making graceful degradation total:
not only does a single failing source not abort the call, but a call in which \emph{every}
source fails still returns a structured result instead of an error. We \emph{loosened} the
optional persistence choice: rather than PostgreSQL we used a filesystem-JSON cache, which is
persistent, inspectable, and adds no service to the container --- an intentional
simplicity-over-power trade-off for a pure key$\rightarrow$blob workload.

\subsection{Provider choice and the AI module contract}

The LLM is selected at runtime through the provided factory; for development the default is
\texttt{LLM\_PROVIDER=anthropic} with model \texttt{claude-sonnet-4-6}, and the web backend
default is \texttt{WEB\_SEARCH\_PROVIDER=duckduckgo} (no key, convenient for grading). Our
code calls exactly four functions of the frozen package --- \texttt{ai.fetch\_wikipedia},
\texttt{ai.fetch\_arxiv}, \texttt{ai.fetch\_web}, and \texttt{ai.synthesize} --- and the
schemas \texttt{Source}, \texttt{Citation}, \texttt{AnswerWithCitations}.
\textbf{We did not modify the public interface of the provided \texttt{ai/} package in any
way.} Against malformed model output we rely on two guards: the provided synthesizer already
drops out-of-range / hallucinated citation indices, and our service layer validates that a
call returned a usable, non-empty result before caching or returning it, treating provider
faults as the transient \texttt{ProviderError} class described in \S\ref{sec:robustness}.

% =====================================================================
\section{Architecture}\label{sec:arch}

\subsection{Module map}

The system is a strict downward-dependency pipeline. The CLI drives the core researcher, which
drives the concurrency orchestrator, which fans out to the service wrapper, which is the
\emph{only} module that touches the frozen \texttt{ai/} package, the cache, and the config.

\begin{lstlisting}[style=plain]
        CLI  (src/cli.py, researcher/__main__.py)
         |  ask "question" [--sources ...] [--no-cache]
         v
   core/researcher.py        validate -> orchestrate -> synthesize -> ResearchResult
         |
         v
 concurrency/orchestrator.py  asyncio.gather(return_exceptions=True)
         |                    + Semaphore(max_concurrency) + per-source asyncio.timeout
         |  fetch_source(origin, query)  x3, one shared httpx.AsyncClient
         v
 services/ai_service.py       retry+backoff | timeout | logging | cache lookup/store
     |        |         |
     v        v         v
 storage/  ===AI BOUNDARY===   config.py
 cache_store   ai/ (FROZEN)    Settings (env-driven)
 (JSON+TTL)    fetch_* / synthesize
\end{lstlisting}

\textbf{Module-by-module.}
\texttt{config.py} reads \texttt{.env} into a typed \texttt{Settings}.
\texttt{models.py} holds our SE-layer schemas.
\texttt{storage/cache\_store.py} is a TTL-aware filesystem JSON cache.
\texttt{services/ai\_service.py} wraps every \texttt{ai.*} call with retry, timeout, logging,
and caching --- the single crossing of the AI boundary.
\texttt{concurrency/orchestrator.py} runs the three fetches concurrently with bounded
parallelism and graceful degradation.
\texttt{core/researcher.py} validates input and assembles the final \texttt{ResearchResult}.
\texttt{cli.py} parses arguments and renders the answer.
Only the arrows into \texttt{ai\_service} cross the AI boundary; swapping the LLM provider
(Anthropic~$\rightarrow$~OpenAI) changes \emph{zero} of our files --- it is one environment
variable, because the provided factory hides the concrete provider behind an ABC.

\subsection{OOP design \& key abstractions}

Our own defensible abstraction is \emph{composition via a dispatch registry} in the service
layer, rather than inheritance. The three source fetchers share one call shape, so instead of
an \texttt{if/elif} ladder or a subclass per source we map origin~$\rightarrow$~coroutine and
inject the shared client and settings:

\begin{lstlisting}[style=python]
_FETCHERS = {
    "wikipedia": fetch_wikipedia,
    "arxiv": fetch_arxiv,
    "web": fetch_web,
}

async def fetch_source(origin, query, *, client=None, settings, use_cache=True):
    fetcher = _FETCHERS.get(origin)
    if fetcher is None:
        raise ValueError(f"unknown origin {origin!r}")
    ...
    return await _fetch_with_retry(fetcher, query, client=client, settings=settings)
\end{lstlisting}

\textbf{Justification.} We chose composition because the fetchers are stateless functions
with an identical signature --- there is no shared state or behaviour to inherit, so a class
hierarchy would add ceremony without value. A registry keeps ``add a source'' to one line and
keeps retry/timeout/caching in exactly one place. Where the problem genuinely called for an
open provider set with per-provider state (API keys, endpoints), the \emph{provided} code
correctly used an ABC (\texttt{WebSearchProvider} with Tavily/Serper/DuckDuckGo subclasses);
we deliberately did not duplicate that pattern where plain functions suffice.

\subsection{Data flow across module boundaries}

We defined three SE-layer models in \texttt{src/models.py}: \texttt{SourceResult} (per-source
outcome: origin, sources, \texttt{ok}, error, \texttt{elapsed}), \texttt{ResearchResult}
(question, answer, per-source results, \texttt{degraded} flag, total elapsed, \texttt{from\_cache}),
and \texttt{CacheEntry} (key, sources, \texttt{created\_at}). We reused the provided
\texttt{Source} / \texttt{AnswerWithCitations} verbatim wherever the data was exactly what
\texttt{ai/} produced, and wrapped them only when we needed to attach SE concerns the provided
schema does not carry --- timing and success/failure for the orchestrator, cache metadata for
the store. \textbf{No naked dictionaries cross module boundaries; every inter-module value is
a typed pydantic model.}

% =====================================================================
\section{Concurrency \& Performance}\label{sec:concurrency}

\subsection{Why this concurrency model}

The workload is I/O-bound --- three independent network round-trips that spend nearly all
their wall-clock time waiting on remote servers, not on CPU. \texttt{asyncio} is therefore the
right primitive: a single thread can keep all three requests in flight, and there is no CPU
work to parallelize that would justify \texttt{ProcessPoolExecutor}, nor any blocking SDK that
would force a \texttt{ThreadPoolExecutor}. Threads would also work for three requests but cost
one OS thread per source and a shared connection pool is cleaner under one event loop. The
bottleneck the design targets is \emph{total network latency}: sequentially it is the sum of
the three sources; concurrently it collapses to the slowest single source. We bound
parallelism with \texttt{asyncio.Semaphore(max\_concurrency)} (default 5, configurable via
\texttt{MAX\_CONCURRENCY}); at 1 the design degenerates to sequential, and at a very high value
we would risk a provider rate-limit (HTTP 429), so a small bound respects politeness while
still overlapping our three-way fan-out. The concurrent design starts winning at $N\ge2$
sources.

\subsection{Sequential-vs-concurrent benchmark}

\texttt{scripts/benchmark.py} runs the same query twice --- once awaiting each source in turn,
once through the orchestrator --- always bypassing the cache so both do identical work.
Measured on a Windows 11 laptop, Python 3.12.4, caches cleared:

\begin{center}\small
\begin{tabular}{lcccc}
\toprule
\textbf{Workload} & $N$ & \textbf{Sequential (s)} & \textbf{Concurrent (s)} & \textbf{Speedup} \\
\midrule
1 query, 3 sources (wiki, arxiv, web) & 3 & 11.43 & 3.64 & 3.1$\times$ \\
\bottomrule
\end{tabular}
\end{center}

Command:
\begin{lstlisting}[style=plain]
python -m scripts.benchmark "What is photosynthesis?" --sources wiki,arxiv,web
\end{lstlisting}

\textbf{Honest interpretation.} We reached almost exactly the theoretical $3\times$ for a
three-source fan-out, which is expected because the three latencies were comparable and the
work is pure waiting. An important caveat: during measurement all three live endpoints were
failing (see \S\ref{sec:robustness}), so the numbers time the concurrent-vs-sequential
handling of the retry-and-timeout path rather than of successful downloads --- but the
comparison is still valid, because sequential pays each source's retry/timeout budget in turn
(sum) while concurrent pays them in parallel (max). The new bottleneck after parallelizing the
fetch is the \emph{LLM synthesis step}, which is a single call that runs after all fetches
complete and is therefore serial tail latency; pushing the curve further would mean streaming
the synthesis or overlapping it with late-arriving sources.

% =====================================================================
\section{Robustness \& Error Handling}\label{sec:robustness}

\subsection{Wrapping the AI module}

Every \texttt{ai.*} call goes through \texttt{ai\_service}, which adds four things the thin
provided layer lacks. \textbf{Retries:} a \texttt{tenacity} policy retries \emph{only}
\texttt{ProviderError} --- the transient class (network errors, 5xx, 429) --- for up to
\texttt{MAX\_RETRIES} (default 3) attempts with \texttt{wait\_exponential} backoff.
Non-transient failures such as a \texttt{ValueError} from empty input are \emph{not} retried,
because they cannot succeed on repeat and retrying them only wastes the budget.
\textbf{Timeouts:} each attempt is wrapped in \texttt{asyncio.timeout(source\_timeout)}
(default 10\,s) as a hard per-call ceiling. \textbf{Structured logging:} the standard
\texttt{logging} module (never \texttt{print}) records origin, truncated query, result count,
and elapsed time at INFO, full payloads at DEBUG, and each retry at WARNING; the level is set
once from \texttt{LOG\_LEVEL}. \textbf{Validation:} input is checked before the call (see
\S\ref{sec:validation}) and the result is checked for usability before it is cached or
returned. Walking one case end-to-end: a provider 429 is raised by \texttt{ai/} as a
\texttt{ProviderError}, caught by the retry predicate, logged at WARNING, waited-out with
exponential backoff, and retried; after the attempt budget is exhausted the original error is
re-raised (\texttt{reraise=True}) and surfaces to the orchestrator, which records it as a
failed \texttt{SourceResult} without cancelling the other sources.

\subsection{Failure-mode analysis (one concrete incident)}

\textbf{Incident: all three live sources unreachable.} We did not have to inject this failure
--- it occurred naturally, which made it a stronger test. Running
\texttt{python -m researcher ask "How does CRISPR-Cas9 work?"} against the live network, every
source failed for a reason inside the frozen \texttt{ai/} layer: arXiv returned HTTP 301 (the
fetcher issues an \texttt{http://} request and does not follow the redirect to
\texttt{https://}), Wikipedia returned HTTP 403 (the request carries no \texttt{User-Agent}
header), and the pinned \texttt{duckduckgo-search} dependency had been renamed to \texttt{ddgs}
upstream and now returns an empty list.

\emph{(i) How it surfaced:} a normal live query. \emph{(ii) What the system did:} each source
was retried three times with backoff, each hit its timeout independently, and the orchestrator
collected three failed \texttt{SourceResult}s via \texttt{return\_exceptions=True}.
\emph{(iii) What the user saw:} a clean message --- ``No sources could be found for this
question\dots'' followed by ``\emph{Note: wikipedia, arxiv, web could not be reached, so this
answer may be incomplete}'' --- and no stack trace. \emph{(iv) What the logs showed:} an
ordered WARNING trail of \texttt{retrying (1/3)}, \texttt{(2/3)}, then per-source
\texttt{FAILED} lines naming the exact HTTP status, which was enough to diagnose all three
root causes without re-running the code. We originally expected at least DuckDuckGo (keyless)
to succeed; the log is what revealed the upstream package rename. Because \texttt{ai/} is
contractually frozen we could not patch the redirect/header/rename at the source, so the
correct engineering response was exactly what the system already did --- degrade cleanly and
report --- which validated the robustness path against real endpoints.

\subsection{Input validation}\label{sec:validation}

The single entry point is the CLI. \texttt{validate\_question} strips the question and rejects
it if it is shorter than \texttt{MIN\_QUESTION\_LEN} (default 3) or longer than
\texttt{MAX\_QUESTION\_LEN} (default 500), raising a \texttt{ValueError} that the CLI surfaces
as a one-line message rather than a traceback. The \texttt{-{}-sources} flag is validated
against the allowed set \texttt{\{wiki, arxiv, web\}}; an unknown origin raises a clear error.
Environment values are validated by \texttt{pydantic-settings} at load time (wrong types fail
fast). A 50{,}000-character query is rejected by the length bound before any network call, so
an adversarial oversized input cannot reach a provider; cache keys are hashed, so a query can
never escape the storage directory via crafted path characters.

\subsection{Rate limit \& politeness}

Politeness has two layers: the \texttt{Semaphore} bounds concurrent in-flight requests so we
never burst more than \texttt{MAX\_CONCURRENCY} (default 5) at once, and the exponential-backoff
retry acts as a sleep-on-429 strategy --- a provider 429 is a \texttt{ProviderError}, so the
backoff naturally spaces out subsequent attempts. The highest burst our test suite produces is
three simultaneous fetches (one per source), well under the bound. We did not observe a real
429 in development because the live endpoints failed earlier in the request path; the strategy
is therefore exercised by unit tests that inject \texttt{ProviderError} rather than by a
production 429.

% =====================================================================
\section{Storage \& Caching}\label{sec:storage}

\subsection{Design}

We did not use SQLite/PostgreSQL. The workload is a pure key$\rightarrow$value blob store
(\texttt{(origin, query)} $\rightarrow$ list of \texttt{Source}s), for which a relational
schema would be over-engineering and would add a service to the container. Instead each entry
is one JSON file under \texttt{CACHE\_DIR} (default \texttt{.cache/}), named by a stable key.
The key is the source-of-truth of identity; the stored \texttt{Source} list and
\texttt{created\_at} timestamp are the cached payload.

\begin{lstlisting}[style=python]
def make_key(origin, query):
    canonical = f"{origin.lower().strip()}|{query.lower().strip()}"
    return f"{origin}_{hashlib.sha256(canonical.encode()).hexdigest()[:16]}"
\end{lstlisting}

The key is \emph{canonicalised} (lower-cased, stripped) so ``\texttt{WHAT IS X?}'' and
``\texttt{what is x}'' collide onto one entry, then hashed so arbitrary query characters
cannot break the filename or escape the directory.

\subsection{Caching policy}

We cache each \texttt{(origin, query)} result for \texttt{CACHE\_TTL} seconds (default 86{,}400
= 24\,h). On read, the entry's age is compared to the TTL and anything older is treated as a
miss and refetched --- this is time-based invalidation. \texttt{ai\_service} checks the cache
before every fetch and writes non-empty results after; \texttt{-{}-no-cache} bypasses both ends.
In a demo run that issues the same query twice, the first call is a miss and the second is a
\textbf{100\% cache hit} (the underlying fetcher is invoked exactly once, verified by a unit
test that counts fetcher calls). The value most likely to go stale is an LLM-shaped result if
the prompt template changes; in production we would add a prompt-version component to the key
so a changed prompt misses the old entries automatically.

% =====================================================================
\section{Testing}\label{sec:testing}

\subsection{Strategy \& coverage}

Every test runs \textbf{offline}: each \texttt{ai.*} call is replaced with a fake via
\texttt{monkeypatch}, so the suite needs no network and no API keys and is safe in CI. Async
code is driven by \texttt{pytest-asyncio} (\texttt{asyncio\_mode = strict}). We chose not to
chase the last few lines --- three uncovered lines in \texttt{config.py} are the
\texttt{\_\_main\_\_} print block, and the two in the orchestrator are a defensive branch that
only fires on a malformed internal result --- because covering them would test framework
plumbing, not our logic.

\begin{center}\small
\begin{tabular}{lc}
\toprule
\textbf{Suite} & \textbf{Coverage} \\
\midrule
\texttt{src/} (our SE layer) & 98\,\% (249 stmts, 6 missed) \\
Provided \texttt{ai/} smoke tests & passing \\
\bottomrule
\end{tabular}
\end{center}

The full suite is \textbf{51 tests passing} in $\sim$2\,s, and the provided
\texttt{tests/test\_ai\_smoke.py} is unmodified and passes.
A \texttt{mypy src} run reports 6 warnings: one is in the frozen \texttt{ai/} package (which we
may not edit), and five are benign false positives in \texttt{ai\_service.py} --- mypy cannot
prove the three registered fetchers share a signature, nor that the retry loop
(\texttt{reraise=True}) always assigns before use. All are documented; none affects runtime,
as the 51 passing tests confirm.

\subsection{Notable tests}

\textbf{(i) Happy path (end-to-end).} \texttt{test\_fetch\_source\_happy\_path} mocks a fetcher
and asserts a well-formed \texttt{Source} flows back through the cache-write path.
\textbf{(ii) Concurrency.} \texttt{test\_orchestrator} verifies that when one of three
gathered fetches raises, \texttt{return\_exceptions=True} still yields three
\texttt{SourceResult}s --- two \texttt{ok}, one failed --- and no exception propagates.
\textbf{(iii) Error path.} \texttt{test\_fetch\_source\_retries\_then\_succeeds} and
\texttt{\dots\_reraises\_after\_exhaustion} pin the retry contract: a \texttt{ProviderError}
that clears on the third attempt succeeds, while a permanent one is re-raised after exactly
\texttt{MAX\_RETRIES} attempts. \texttt{\dots\_times\_out} confirms a slow fetcher raises
\texttt{TimeoutError} and is \emph{not} retried. The test we wish we had time for is a
load-style test asserting the semaphore never lets more than \texttt{MAX\_CONCURRENCY}
coroutines run at once under a large fan-out.

% =====================================================================
\section{Deployment \& Reproducibility}\label{sec:deploy}

\subsection{Docker image}

\begin{lstlisting}[style=plain]
docker build -t research-assistant .
docker run --rm research-assistant        # runs: python demo_ai.py --offline
\end{lstlisting}

It is a \textbf{single-stage} build on \texttt{python:3.12-slim}. The image is \textbf{76.5\,MB
content} (324\,MB on-disk including the shared base layers), comparable to a fresh
\texttt{python:3.12-slim} ($\sim$120\,MB) because our only added weight is the pinned Python
dependencies and the small source tree; a \texttt{.dockerignore} keeps the virtualenv, caches,
git history, and \texttt{.env} out of the image. Dependencies are installed before the source
is copied so the dependency layer is cached across code changes. The default command runs the
offline demo end-to-end with no keys or network; the live CLI can be run by passing env vars
with \texttt{-e} / \texttt{-{}-env-file}.

\subsection{Configuration}

Every variable the app reads, and the effect if it is missing:
\texttt{LLM\_PROVIDER} (default \texttt{anthropic}); the provider's API key
(\texttt{ANTHROPIC\_API\_KEY} etc.) --- if missing, live synthesis raises a clear
\texttt{ProviderError}, but the offline demo and tests are unaffected;
\texttt{WEB\_SEARCH\_PROVIDER} (default \texttt{duckduckgo}); \texttt{TAVILY\_API\_KEY} /
\texttt{SERPER\_API\_KEY} (needed only for those backends); \texttt{CACHE\_DIR} (default
\texttt{.cache}); \texttt{CACHE\_TTL} (default 86400); \texttt{SOURCE\_TIMEOUT} (default 10);
\texttt{MAX\_CONCURRENCY} (default 5); \texttt{MAX\_RETRIES} / \texttt{RETRY\_BACKOFF}
(defaults 3 / 1.0); \texttt{MIN\_QUESTION\_LEN} / \texttt{MAX\_QUESTION\_LEN} (defaults 3 /
500); \texttt{LOG\_LEVEL} (default \texttt{INFO}). Every non-key variable has a safe default,
so the app runs with an empty \texttt{.env}.

% =====================================================================
\section{Limitations \& Future Work}\label{sec:limits}

\begin{itemize}[leftmargin=1.4em,itemsep=3pt]
  \item \textbf{Live retrieval depends on frozen code we cannot patch.} arXiv (301 redirect),
  Wikipedia (403 / missing User-Agent), and the renamed \texttt{duckduckgo-search} package all
  fail inside \texttt{ai/}; we validate the pipeline via the offline demo and mocked tests. A
  permitted fix would be a thin request-shim layer of our own that adds headers and follows
  redirects \emph{before} handing off, without editing \texttt{ai/}.
  \item \textbf{The filesystem cache has no concurrent-writer locking} and grows one file per
  unique query; a production version would add locking or move to SQLite/Redis and an eviction
  policy.
  \item \textbf{Synthesis is serial tail latency} --- the single LLM call runs after all
  fetches, so it dominates once retrieval is parallel; streaming would help.
  \item \textbf{Rate limiting is request-count only}, not token-aware; a busy day with a real
  provider could still approach a tokens-per-minute limit the semaphore does not see.
  \item \textbf{Two-person scope} meant we prioritized the required layer over optional bonuses
  (multi-provider failover, cost telemetry, CI); these are the natural next week of work.
\end{itemize}

% =====================================================================
\section{Tools \& Acknowledgements}\label{sec:tools}

We used AI coding assistants as collaborators, and disclose this specifically rather than
generically. \textbf{Claude} was used to scaffold the SE layer module by module and to explain
each concept (\texttt{asyncio}, \texttt{tenacity} retry semantics, \texttt{pytest} mocking,
Docker) as we built; every module was then reviewed, run, and unit-tested by its owner, and we
can each defend our own files. Concrete example: the \texttt{ai\_service} retry policy was
drafted with Claude, but we tuned it after reading the live logs --- narrowing the retry
predicate to \texttt{ProviderError} only, so that non-transient \texttt{ValueError}s are not
uselessly retried --- and we added the per-attempt \texttt{asyncio.timeout} after observing
that a hung source would otherwise block within the retry budget. The provided \texttt{ai/}
package, its schemas, and \texttt{demo\_ai.py} are course-supplied and were left unmodified.

% =====================================================================
\section*{References}
\addcontentsline{toc}{section}{References}

\begin{enumerate}[label={[\arabic*]},leftmargin=2.4em,itemsep=2pt]
  \item Anthropic API documentation. \url{https://docs.anthropic.com/}
  \item \texttt{tenacity} --- retrying library. \url{https://tenacity.readthedocs.io/}
  \item \texttt{pydantic-settings} documentation. \url{https://docs.pydantic.dev/latest/concepts/pydantic_settings/}
  \item \texttt{httpx} --- async HTTP client. \url{https://www.python-httpx.org/}
  \item \texttt{pytest-asyncio} documentation. \url{https://pytest-asyncio.readthedocs.io/}
  \item Python \texttt{asyncio} --- \texttt{gather}, \texttt{timeout}, \texttt{Semaphore}. \url{https://docs.python.org/3/library/asyncio.html}
\end{enumerate}

\appendix
% =====================================================================
\section{Appendix --- Reproducing the benchmark}
Anyone with the repo and a Python 3.11+ environment can reproduce the \S\ref{sec:concurrency}
table:

\begin{lstlisting}[style=plain]
git clone https://github.com/Yunis-Allahverdi/async-research-assistant
cd async-research-assistant
python -m venv venv && source venv/bin/activate   # Windows: venv\Scripts\Activate.ps1
python -m pip install -r requirements-ai.txt -r requirements.txt
python -m scripts.benchmark "What is photosynthesis?" --sources wiki,arxiv,web
\end{lstlisting}

\vspace{1em}\hrule\vspace{0.6em}
\noindent\small\textit{End of report --- thank you for reading.}

\end{document}% =====================================================================
% AI Academy -- National AI Center
% Software Engineering Final Project -- Team Report
% Topic 4 -- Async Research Assistant
% =====================================================================

\documentclass[11pt,a4paper]{article}

% ===== EARLY MACROS ==================================================
\usepackage[dvipsnames]{xcolor}
\definecolor{accent2}{HTML}{8B0000}
\definecolor{ph}{HTML}{6B7280}

% ===== CONFIGURATION =================================================
\newcommand{\teamname}{Allahverdi \& Isgandarli}
\newcommand{\topicchoice}{Topic 4 -- Async Research Assistant}
\newcommand{\member}[2]{\textbf{#1} (\texttt{#2})}
\newcommand{\reportdate}{18 September 2026}
\newcommand{\repolink}{\url{https://github.com/Yunis-Allahverdi/async-research-assistant}}

% ===== course meta ==================================================
\newcommand{\coursecode}{AI-ENG-110}
\newcommand{\coursename}{Software Engineering}
\newcommand{\institution}{AI Academy, National AI Center}
\newcommand{\semester}{Spring 2026}

% ===== PACKAGES =====================================================
\usepackage[margin=2.2cm,headheight=15pt]{geometry}
\usepackage[T1]{fontenc}
\usepackage[utf8]{inputenc}
\usepackage[final,protrusion=true,expansion=false]{microtype}
\usepackage{amsmath,amssymb}
\usepackage{graphicx}
\usepackage{booktabs}
\usepackage{array}
\usepackage{tabularx}
\usepackage{ragged2e}
\usepackage{enumitem}
\usepackage[parfill]{parskip}
\usepackage{listings}
\usepackage[most]{tcolorbox}
\usepackage{titlesec}
\usepackage{fancyhdr}
\usepackage{lastpage}
\usepackage[hidelinks]{hyperref}

\emergencystretch=3em
\hbadness=10000
\hfuzz=2pt

\hypersetup{
  colorlinks=true,
  linkcolor=NavyBlue,
  urlcolor=NavyBlue,
  citecolor=NavyBlue,
  pdftitle={\coursecode\ Final Project Report},
  pdfauthor={\teamname}
}

% ===== STYLING ======================================================
\definecolor{accent}{HTML}{0B3D91}
\definecolor{soft}{HTML}{F2F4F8}

\titleformat{\section}
  {\Large\bfseries\color{accent}}{\thesection.}{0.6em}{}
\titleformat{\subsection}
  {\large\bfseries\color{accent}}{\thesubsection}{0.5em}{}

\definecolor{codegreen}{rgb}{0,0.55,0}
\definecolor{codegray}{rgb}{0.4,0.4,0.4}
\definecolor{codepurple}{rgb}{0.55,0,0.78}
\definecolor{backcolour}{rgb}{0.97,0.97,0.97}

\lstdefinestyle{python}{
  language=Python,
  backgroundcolor=\color{backcolour},
  commentstyle=\color{codegreen},
  keywordstyle=\color{accent},
  numberstyle=\tiny\color{codegray},
  stringstyle=\color{codepurple},
  basicstyle=\ttfamily\footnotesize,
  breakatwhitespace=false,
  breaklines=true,
  keepspaces=true,
  numbers=left,
  numbersep=6pt,
  tabsize=2,
  frame=single,
  framerule=0.4pt,
  rulecolor=\color{black!30},
}
\lstdefinestyle{plain}{
  backgroundcolor=\color{backcolour},
  basicstyle=\ttfamily\footnotesize,
  breaklines=true,
  numbers=none,
  frame=single,
  framerule=0.4pt,
  rulecolor=\color{black!30},
}

\pagestyle{fancy}
\fancyhf{}
\fancyhead[L]{\small\textit{\coursecode\ -- Final Report -- \teamname}}
\fancyhead[R]{\small\textit{\semester}}
\fancyfoot[L]{\small\textit{\institution}}
\fancyfoot[R]{\small Page \thepage\ of \pageref{LastPage}}
\renewcommand{\headrulewidth}{0.4pt}
\renewcommand{\footrulewidth}{0.4pt}

% =====================================================================
\begin{document}

\begin{center}
{\small \textsc{\institution} \\[0.2em] \textsc{\coursecode\ -- \coursename\ -- \semester}}\\[1.2em]
{\Large\bfseries\color{accent} Final Project Report}\\[0.3em]
{\LARGE\bfseries \topicchoice}\\[0.6em]
\rule{0.85\textwidth}{0.6pt}
\end{center}

\vspace{0.6em}
\begin{tabularx}{\textwidth}{@{}lXlX@{}}
\textbf{Team:} & \teamname & \textbf{Date:} & \reportdate \\
\textbf{Repo:} & \repolink & \textbf{Tag:} & \texttt{v1.0-final} \\
\textbf{Members:} & \multicolumn{3}{X@{}}{%
  \member{Yunis Allahverdi}{@Yunis-Allahverdi},\quad
  \member{Fariz Isgandarli}{@farizisgenderli05-create}
}\\
\end{tabularx}
\vspace{0.4em}\hrule\vspace{0.6em}

% ===== EXECUTIVE SUMMARY =============================================
\section*{Executive Summary}
\addcontentsline{toc}{section}{Executive Summary}

We built an asynchronous research assistant: given a natural-language question, it queries
Wikipedia, arXiv, and a web-search provider concurrently, then uses a provider-agnostic LLM
to synthesize a single answer carrying inline \texttt{[N]} citations back to the retrieved
sources. The provided \texttt{ai/} package (fetchers plus citation synthesizer) was treated
as a frozen dependency; our work was the entire software-engineering layer around it ---
typed configuration, an \texttt{asyncio} orchestration layer, a filesystem cache, a retry /
timeout / logging service wrapper, input validation, a CLI, an offline test suite, and a
Docker image. The one number that proves the concurrency design works: on an identical
three-source workload our benchmark drops wall-clock time from \textbf{11.43\,s sequential to
3.64\,s concurrent (3.1$\times$)}. Our headline robustness story is that when \emph{all three}
live sources failed for reasons inside the frozen \texttt{ai/} layer, the system retried with
backoff, hit its per-source timeouts independently, and degraded to a clean ``no sources
found'' response rather than crashing. The most surprising lesson was that the hardest part
of an ``AI project'' was not the AI at all but everything that makes a thin API robust,
observable, and reproducible; if we started over we would push the retry policy to be
token-aware and add a real HTTP-header shim layer we are permitted to own.

% ===== TABLE OF CONTENTS ============================================
\tableofcontents
\vspace{0.6em}\hrule\vspace{0.4em}

% =====================================================================
\section{Project Overview}

\subsection{Problem framing \& scope}

The system answers a research question by gathering short, attributed excerpts from three
independent sources and letting a language model fuse them into one concise, cited answer.
The \emph{input} is a single question string; the \emph{outputs} are a synthesized answer
plus a numbered reference list, where every \texttt{[N]} marker in the prose maps to one
retrieved source. The only user-facing entry point is a command-line interface:
\texttt{python -m researcher ask "\dots"}, with \texttt{-{}-sources} to restrict the source
set and \texttt{-{}-no-cache} to bypass the cache. No HTTP server is required for this topic.

The provider stack is deliberately pluggable and selected entirely by environment variables:
the LLM is provider-agnostic (Anthropic / OpenAI / Gemini via \texttt{LLM\_PROVIDER}), and the
web-search backend is one of Tavily / Serper / DuckDuckGo via \texttt{WEB\_SEARCH\_PROVIDER}.
Wikipedia and arXiv need no key.

Relative to \texttt{TOPIC.md}, we implemented the full required SE layer and made two
deliberate scope calls. We \emph{tightened} the spec by making graceful degradation total:
not only does a single failing source not abort the call, but a call in which \emph{every}
source fails still returns a structured result instead of an error. We \emph{loosened} the
optional persistence choice: rather than PostgreSQL we used a filesystem-JSON cache, which is
persistent, inspectable, and adds no service to the container --- an intentional
simplicity-over-power trade-off for a pure key$\rightarrow$blob workload.

\subsection{Provider choice and the AI module contract}

The LLM is selected at runtime through the provided factory; for development the default is
\texttt{LLM\_PROVIDER=anthropic} with model \texttt{claude-sonnet-4-6}, and the web backend
default is \texttt{WEB\_SEARCH\_PROVIDER=duckduckgo} (no key, convenient for grading). Our
code calls exactly four functions of the frozen package --- \texttt{ai.fetch\_wikipedia},
\texttt{ai.fetch\_arxiv}, \texttt{ai.fetch\_web}, and \texttt{ai.synthesize} --- and the
schemas \texttt{Source}, \texttt{Citation}, \texttt{AnswerWithCitations}.
\textbf{We did not modify the public interface of the provided \texttt{ai/} package in any
way.} Against malformed model output we rely on two guards: the provided synthesizer already
drops out-of-range / hallucinated citation indices, and our service layer validates that a
call returned a usable, non-empty result before caching or returning it, treating provider
faults as the transient \texttt{ProviderError} class described in \S\ref{sec:robustness}.

% =====================================================================
\section{Architecture}\label{sec:arch}

\subsection{Module map}

The system is a strict downward-dependency pipeline. The CLI drives the core researcher, which
drives the concurrency orchestrator, which fans out to the service wrapper, which is the
\emph{only} module that touches the frozen \texttt{ai/} package, the cache, and the config.

\begin{lstlisting}[style=plain]
        CLI  (src/cli.py, researcher/__main__.py)
         |  ask "question" [--sources ...] [--no-cache]
         v
   core/researcher.py        validate -> orchestrate -> synthesize -> ResearchResult
         |
         v
 concurrency/orchestrator.py  asyncio.gather(return_exceptions=True)
         |                    + Semaphore(max_concurrency) + per-source asyncio.timeout
         |  fetch_source(origin, query)  x3, one shared httpx.AsyncClient
         v
 services/ai_service.py       retry+backoff | timeout | logging | cache lookup/store
     |        |         |
     v        v         v
 storage/  ===AI BOUNDARY===   config.py
 cache_store   ai/ (FROZEN)    Settings (env-driven)
 (JSON+TTL)    fetch_* / synthesize
\end{lstlisting}

\textbf{Module-by-module.}
\texttt{config.py} reads \texttt{.env} into a typed \texttt{Settings}.
\texttt{models.py} holds our SE-layer schemas.
\texttt{storage/cache\_store.py} is a TTL-aware filesystem JSON cache.
\texttt{services/ai\_service.py} wraps every \texttt{ai.*} call with retry, timeout, logging,
and caching --- the single crossing of the AI boundary.
\texttt{concurrency/orchestrator.py} runs the three fetches concurrently with bounded
parallelism and graceful degradation.
\texttt{core/researcher.py} validates input and assembles the final \texttt{ResearchResult}.
\texttt{cli.py} parses arguments and renders the answer.
Only the arrows into \texttt{ai\_service} cross the AI boundary; swapping the LLM provider
(Anthropic~$\rightarrow$~OpenAI) changes \emph{zero} of our files --- it is one environment
variable, because the provided factory hides the concrete provider behind an ABC.

\subsection{OOP design \& key abstractions}

Our own defensible abstraction is \emph{composition via a dispatch registry} in the service
layer, rather than inheritance. The three source fetchers share one call shape, so instead of
an \texttt{if/elif} ladder or a subclass per source we map origin~$\rightarrow$~coroutine and
inject the shared client and settings:

\begin{lstlisting}[style=python]
_FETCHERS = {
    "wikipedia": fetch_wikipedia,
    "arxiv": fetch_arxiv,
    "web": fetch_web,
}

async def fetch_source(origin, query, *, client=None, settings, use_cache=True):
    fetcher = _FETCHERS.get(origin)
    if fetcher is None:
        raise ValueError(f"unknown origin {origin!r}")
    ...
    return await _fetch_with_retry(fetcher, query, client=client, settings=settings)
\end{lstlisting}

\textbf{Justification.} We chose composition because the fetchers are stateless functions
with an identical signature --- there is no shared state or behaviour to inherit, so a class
hierarchy would add ceremony without value. A registry keeps ``add a source'' to one line and
keeps retry/timeout/caching in exactly one place. Where the problem genuinely called for an
open provider set with per-provider state (API keys, endpoints), the \emph{provided} code
correctly used an ABC (\texttt{WebSearchProvider} with Tavily/Serper/DuckDuckGo subclasses);
we deliberately did not duplicate that pattern where plain functions suffice.

\subsection{Data flow across module boundaries}

We defined three SE-layer models in \texttt{src/models.py}: \texttt{SourceResult} (per-source
outcome: origin, sources, \texttt{ok}, error, \texttt{elapsed}), \texttt{ResearchResult}
(question, answer, per-source results, \texttt{degraded} flag, total elapsed, \texttt{from\_cache}),
and \texttt{CacheEntry} (key, sources, \texttt{created\_at}). We reused the provided
\texttt{Source} / \texttt{AnswerWithCitations} verbatim wherever the data was exactly what
\texttt{ai/} produced, and wrapped them only when we needed to attach SE concerns the provided
schema does not carry --- timing and success/failure for the orchestrator, cache metadata for
the store. \textbf{No naked dictionaries cross module boundaries; every inter-module value is
a typed pydantic model.}

% =====================================================================
\section{Concurrency \& Performance}\label{sec:concurrency}

\subsection{Why this concurrency model}

The workload is I/O-bound --- three independent network round-trips that spend nearly all
their wall-clock time waiting on remote servers, not on CPU. \texttt{asyncio} is therefore the
right primitive: a single thread can keep all three requests in flight, and there is no CPU
work to parallelize that would justify \texttt{ProcessPoolExecutor}, nor any blocking SDK that
would force a \texttt{ThreadPoolExecutor}. Threads would also work for three requests but cost
one OS thread per source and a shared connection pool is cleaner under one event loop. The
bottleneck the design targets is \emph{total network latency}: sequentially it is the sum of
the three sources; concurrently it collapses to the slowest single source. We bound
parallelism with \texttt{asyncio.Semaphore(max\_concurrency)} (default 5, configurable via
\texttt{MAX\_CONCURRENCY}); at 1 the design degenerates to sequential, and at a very high value
we would risk a provider rate-limit (HTTP 429), so a small bound respects politeness while
still overlapping our three-way fan-out. The concurrent design starts winning at $N\ge2$
sources.

\subsection{Sequential-vs-concurrent benchmark}

\texttt{scripts/benchmark.py} runs the same query twice --- once awaiting each source in turn,
once through the orchestrator --- always bypassing the cache so both do identical work.
Measured on a Windows 11 laptop, Python 3.12.4, caches cleared:

\begin{center}\small
\begin{tabular}{lcccc}
\toprule
\textbf{Workload} & $N$ & \textbf{Sequential (s)} & \textbf{Concurrent (s)} & \textbf{Speedup} \\
\midrule
1 query, 3 sources (wiki, arxiv, web) & 3 & 11.43 & 3.64 & 3.1$\times$ \\
\bottomrule
\end{tabular}
\end{center}

Command:
\begin{lstlisting}[style=plain]
python -m scripts.benchmark "What is photosynthesis?" --sources wiki,arxiv,web
\end{lstlisting}

\textbf{Honest interpretation.} We reached almost exactly the theoretical $3\times$ for a
three-source fan-out, which is expected because the three latencies were comparable and the
work is pure waiting. An important caveat: during measurement all three live endpoints were
failing (see \S\ref{sec:robustness}), so the numbers time the concurrent-vs-sequential
handling of the retry-and-timeout path rather than of successful downloads --- but the
comparison is still valid, because sequential pays each source's retry/timeout budget in turn
(sum) while concurrent pays them in parallel (max). The new bottleneck after parallelizing the
fetch is the \emph{LLM synthesis step}, which is a single call that runs after all fetches
complete and is therefore serial tail latency; pushing the curve further would mean streaming
the synthesis or overlapping it with late-arriving sources.

% =====================================================================
\section{Robustness \& Error Handling}\label{sec:robustness}

\subsection{Wrapping the AI module}

Every \texttt{ai.*} call goes through \texttt{ai\_service}, which adds four things the thin
provided layer lacks. \textbf{Retries:} a \texttt{tenacity} policy retries \emph{only}
\texttt{ProviderError} --- the transient class (network errors, 5xx, 429) --- for up to
\texttt{MAX\_RETRIES} (default 3) attempts with \texttt{wait\_exponential} backoff.
Non-transient failures such as a \texttt{ValueError} from empty input are \emph{not} retried,
because they cannot succeed on repeat and retrying them only wastes the budget.
\textbf{Timeouts:} each attempt is wrapped in \texttt{asyncio.timeout(source\_timeout)}
(default 10\,s) as a hard per-call ceiling. \textbf{Structured logging:} the standard
\texttt{logging} module (never \texttt{print}) records origin, truncated query, result count,
and elapsed time at INFO, full payloads at DEBUG, and each retry at WARNING; the level is set
once from \texttt{LOG\_LEVEL}. \textbf{Validation:} input is checked before the call (see
\S\ref{sec:validation}) and the result is checked for usability before it is cached or
returned. Walking one case end-to-end: a provider 429 is raised by \texttt{ai/} as a
\texttt{ProviderError}, caught by the retry predicate, logged at WARNING, waited-out with
exponential backoff, and retried; after the attempt budget is exhausted the original error is
re-raised (\texttt{reraise=True}) and surfaces to the orchestrator, which records it as a
failed \texttt{SourceResult} without cancelling the other sources.

\subsection{Failure-mode analysis (one concrete incident)}

\textbf{Incident: all three live sources unreachable.} We did not have to inject this failure
--- it occurred naturally, which made it a stronger test. Running
\texttt{python -m researcher ask "How does CRISPR-Cas9 work?"} against the live network, every
source failed for a reason inside the frozen \texttt{ai/} layer: arXiv returned HTTP 301 (the
fetcher issues an \texttt{http://} request and does not follow the redirect to
\texttt{https://}), Wikipedia returned HTTP 403 (the request carries no \texttt{User-Agent}
header), and the pinned \texttt{duckduckgo-search} dependency had been renamed to \texttt{ddgs}
upstream and now returns an empty list.

\emph{(i) How it surfaced:} a normal live query. \emph{(ii) What the system did:} each source
was retried three times with backoff, each hit its timeout independently, and the orchestrator
collected three failed \texttt{SourceResult}s via \texttt{return\_exceptions=True}.
\emph{(iii) What the user saw:} a clean message --- ``No sources could be found for this
question\dots'' followed by ``\emph{Note: wikipedia, arxiv, web could not be reached, so this
answer may be incomplete}'' --- and no stack trace. \emph{(iv) What the logs showed:} an
ordered WARNING trail of \texttt{retrying (1/3)}, \texttt{(2/3)}, then per-source
\texttt{FAILED} lines naming the exact HTTP status, which was enough to diagnose all three
root causes without re-running the code. We originally expected at least DuckDuckGo (keyless)
to succeed; the log is what revealed the upstream package rename. Because \texttt{ai/} is
contractually frozen we could not patch the redirect/header/rename at the source, so the
correct engineering response was exactly what the system already did --- degrade cleanly and
report --- which validated the robustness path against real endpoints.

\subsection{Input validation}\label{sec:validation}

The single entry point is the CLI. \texttt{validate\_question} strips the question and rejects
it if it is shorter than \texttt{MIN\_QUESTION\_LEN} (default 3) or longer than
\texttt{MAX\_QUESTION\_LEN} (default 500), raising a \texttt{ValueError} that the CLI surfaces
as a one-line message rather than a traceback. The \texttt{-{}-sources} flag is validated
against the allowed set \texttt{\{wiki, arxiv, web\}}; an unknown origin raises a clear error.
Environment values are validated by \texttt{pydantic-settings} at load time (wrong types fail
fast). A 50{,}000-character query is rejected by the length bound before any network call, so
an adversarial oversized input cannot reach a provider; cache keys are hashed, so a query can
never escape the storage directory via crafted path characters.

\subsection{Rate limit \& politeness}

Politeness has two layers: the \texttt{Semaphore} bounds concurrent in-flight requests so we
never burst more than \texttt{MAX\_CONCURRENCY} (default 5) at once, and the exponential-backoff
retry acts as a sleep-on-429 strategy --- a provider 429 is a \texttt{ProviderError}, so the
backoff naturally spaces out subsequent attempts. The highest burst our test suite produces is
three simultaneous fetches (one per source), well under the bound. We did not observe a real
429 in development because the live endpoints failed earlier in the request path; the strategy
is therefore exercised by unit tests that inject \texttt{ProviderError} rather than by a
production 429.

% =====================================================================
\section{Storage \& Caching}\label{sec:storage}

\subsection{Design}

We did not use SQLite/PostgreSQL. The workload is a pure key$\rightarrow$value blob store
(\texttt{(origin, query)} $\rightarrow$ list of \texttt{Source}s), for which a relational
schema would be over-engineering and would add a service to the container. Instead each entry
is one JSON file under \texttt{CACHE\_DIR} (default \texttt{.cache/}), named by a stable key.
The key is the source-of-truth of identity; the stored \texttt{Source} list and
\texttt{created\_at} timestamp are the cached payload.

\begin{lstlisting}[style=python]
def make_key(origin, query):
    canonical = f"{origin.lower().strip()}|{query.lower().strip()}"
    return f"{origin}_{hashlib.sha256(canonical.encode()).hexdigest()[:16]}"
\end{lstlisting}

The key is \emph{canonicalised} (lower-cased, stripped) so ``\texttt{WHAT IS X?}'' and
``\texttt{what is x}'' collide onto one entry, then hashed so arbitrary query characters
cannot break the filename or escape the directory.

\subsection{Caching policy}

We cache each \texttt{(origin, query)} result for \texttt{CACHE\_TTL} seconds (default 86{,}400
= 24\,h). On read, the entry's age is compared to the TTL and anything older is treated as a
miss and refetched --- this is time-based invalidation. \texttt{ai\_service} checks the cache
before every fetch and writes non-empty results after; \texttt{-{}-no-cache} bypasses both ends.
In a demo run that issues the same query twice, the first call is a miss and the second is a
\textbf{100\% cache hit} (the underlying fetcher is invoked exactly once, verified by a unit
test that counts fetcher calls). The value most likely to go stale is an LLM-shaped result if
the prompt template changes; in production we would add a prompt-version component to the key
so a changed prompt misses the old entries automatically.

% =====================================================================
\section{Testing}\label{sec:testing}

\subsection{Strategy \& coverage}

Every test runs \textbf{offline}: each \texttt{ai.*} call is replaced with a fake via
\texttt{monkeypatch}, so the suite needs no network and no API keys and is safe in CI. Async
code is driven by \texttt{pytest-asyncio} (\texttt{asyncio\_mode = strict}). We chose not to
chase the last few lines --- three uncovered lines in \texttt{config.py} are the
\texttt{\_\_main\_\_} print block, and the two in the orchestrator are a defensive branch that
only fires on a malformed internal result --- because covering them would test framework
plumbing, not our logic.

\begin{center}\small
\begin{tabular}{lc}
\toprule
\textbf{Suite} & \textbf{Coverage} \\
\midrule
\texttt{src/} (our SE layer) & 98\,\% (249 stmts, 6 missed) \\
Provided \texttt{ai/} smoke tests & passing \\
\bottomrule
\end{tabular}
\end{center}

The full suite is \textbf{51 tests passing} in $\sim$2\,s, and the provided
\texttt{tests/test\_ai\_smoke.py} is unmodified and passes.
A \texttt{mypy src} run reports 6 warnings: one is in the frozen \texttt{ai/} package (which we
may not edit), and five are benign false positives in \texttt{ai\_service.py} --- mypy cannot
prove the three registered fetchers share a signature, nor that the retry loop
(\texttt{reraise=True}) always assigns before use. All are documented; none affects runtime,
as the 51 passing tests confirm.

\subsection{Notable tests}

\textbf{(i) Happy path (end-to-end).} \texttt{test\_fetch\_source\_happy\_path} mocks a fetcher
and asserts a well-formed \texttt{Source} flows back through the cache-write path.
\textbf{(ii) Concurrency.} \texttt{test\_orchestrator} verifies that when one of three
gathered fetches raises, \texttt{return\_exceptions=True} still yields three
\texttt{SourceResult}s --- two \texttt{ok}, one failed --- and no exception propagates.
\textbf{(iii) Error path.} \texttt{test\_fetch\_source\_retries\_then\_succeeds} and
\texttt{\dots\_reraises\_after\_exhaustion} pin the retry contract: a \texttt{ProviderError}
that clears on the third attempt succeeds, while a permanent one is re-raised after exactly
\texttt{MAX\_RETRIES} attempts. \texttt{\dots\_times\_out} confirms a slow fetcher raises
\texttt{TimeoutError} and is \emph{not} retried. The test we wish we had time for is a
load-style test asserting the semaphore never lets more than \texttt{MAX\_CONCURRENCY}
coroutines run at once under a large fan-out.

% =====================================================================
\section{Deployment \& Reproducibility}\label{sec:deploy}

\subsection{Docker image}

\begin{lstlisting}[style=plain]
docker build -t research-assistant .
docker run --rm research-assistant        # runs: python demo_ai.py --offline
\end{lstlisting}

It is a \textbf{single-stage} build on \texttt{python:3.12-slim}. The image is \textbf{76.5\,MB
content} (324\,MB on-disk including the shared base layers), comparable to a fresh
\texttt{python:3.12-slim} ($\sim$120\,MB) because our only added weight is the pinned Python
dependencies and the small source tree; a \texttt{.dockerignore} keeps the virtualenv, caches,
git history, and \texttt{.env} out of the image. Dependencies are installed before the source
is copied so the dependency layer is cached across code changes. The default command runs the
offline demo end-to-end with no keys or network; the live CLI can be run by passing env vars
with \texttt{-e} / \texttt{-{}-env-file}.

\subsection{Configuration}

Every variable the app reads, and the effect if it is missing:
\texttt{LLM\_PROVIDER} (default \texttt{anthropic}); the provider's API key
(\texttt{ANTHROPIC\_API\_KEY} etc.) --- if missing, live synthesis raises a clear
\texttt{ProviderError}, but the offline demo and tests are unaffected;
\texttt{WEB\_SEARCH\_PROVIDER} (default \texttt{duckduckgo}); \texttt{TAVILY\_API\_KEY} /
\texttt{SERPER\_API\_KEY} (needed only for those backends); \texttt{CACHE\_DIR} (default
\texttt{.cache}); \texttt{CACHE\_TTL} (default 86400); \texttt{SOURCE\_TIMEOUT} (default 10);
\texttt{MAX\_CONCURRENCY} (default 5); \texttt{MAX\_RETRIES} / \texttt{RETRY\_BACKOFF}
(defaults 3 / 1.0); \texttt{MIN\_QUESTION\_LEN} / \texttt{MAX\_QUESTION\_LEN} (defaults 3 /
500); \texttt{LOG\_LEVEL} (default \texttt{INFO}). Every non-key variable has a safe default,
so the app runs with an empty \texttt{.env}.

% =====================================================================
\section{Limitations \& Future Work}\label{sec:limits}

\begin{itemize}[leftmargin=1.4em,itemsep=3pt]
  \item \textbf{Live retrieval depends on frozen code we cannot patch.} arXiv (301 redirect),
  Wikipedia (403 / missing User-Agent), and the renamed \texttt{duckduckgo-search} package all
  fail inside \texttt{ai/}; we validate the pipeline via the offline demo and mocked tests. A
  permitted fix would be a thin request-shim layer of our own that adds headers and follows
  redirects \emph{before} handing off, without editing \texttt{ai/}.
  \item \textbf{The filesystem cache has no concurrent-writer locking} and grows one file per
  unique query; a production version would add locking or move to SQLite/Redis and an eviction
  policy.
  \item \textbf{Synthesis is serial tail latency} --- the single LLM call runs after all
  fetches, so it dominates once retrieval is parallel; streaming would help.
  \item \textbf{Rate limiting is request-count only}, not token-aware; a busy day with a real
  provider could still approach a tokens-per-minute limit the semaphore does not see.
  \item \textbf{Two-person scope} meant we prioritized the required layer over optional bonuses
  (multi-provider failover, cost telemetry, CI); these are the natural next week of work.
\end{itemize}

% =====================================================================
\section{Tools \& Acknowledgements}\label{sec:tools}

We used AI coding assistants as collaborators, and disclose this specifically rather than
generically. \textbf{Claude} was used to scaffold the SE layer module by module and to explain
each concept (\texttt{asyncio}, \texttt{tenacity} retry semantics, \texttt{pytest} mocking,
Docker) as we built; every module was then reviewed, run, and unit-tested by its owner, and we
can each defend our own files. Concrete example: the \texttt{ai\_service} retry policy was
drafted with Claude, but we tuned it after reading the live logs --- narrowing the retry
predicate to \texttt{ProviderError} only, so that non-transient \texttt{ValueError}s are not
uselessly retried --- and we added the per-attempt \texttt{asyncio.timeout} after observing
that a hung source would otherwise block within the retry budget. The provided \texttt{ai/}
package, its schemas, and \texttt{demo\_ai.py} are course-supplied and were left unmodified.

% =====================================================================
\section*{References}
\addcontentsline{toc}{section}{References}

\begin{enumerate}[label={[\arabic*]},leftmargin=2.4em,itemsep=2pt]
  \item Anthropic API documentation. \url{https://docs.anthropic.com/}
  \item \texttt{tenacity} --- retrying library. \url{https://tenacity.readthedocs.io/}
  \item \texttt{pydantic-settings} documentation. \url{https://docs.pydantic.dev/latest/concepts/pydantic_settings/}
  \item \texttt{httpx} --- async HTTP client. \url{https://www.python-httpx.org/}
  \item \texttt{pytest-asyncio} documentation. \url{https://pytest-asyncio.readthedocs.io/}
  \item Python \texttt{asyncio} --- \texttt{gather}, \texttt{timeout}, \texttt{Semaphore}. \url{https://docs.python.org/3/library/asyncio.html}
\end{enumerate}

\appendix
% =====================================================================
\section{Appendix --- Reproducing the benchmark}
Anyone with the repo and a Python 3.11+ environment can reproduce the \S\ref{sec:concurrency}
table:

\begin{lstlisting}[style=plain]
git clone https://github.com/Yunis-Allahverdi/async-research-assistant
cd async-research-assistant
python -m venv venv && source venv/bin/activate   # Windows: venv\Scripts\Activate.ps1
python -m pip install -r requirements-ai.txt -r requirements.txt
python -m scripts.benchmark "What is photosynthesis?" --sources wiki,arxiv,web
\end{lstlisting}

\vspace{1em}\hrule\vspace{0.6em}
\noindent\small\textit{End of report --- thank you for reading.}

\end{document}